\section{Signed tabloids and conditional products}
\label{l:sec:tabloid}

Information about the images of almost all positions can control a
representation through a larger space with a particularly simple basis.
In a signed tabloid representation, each permutation permutes the basis
vectors and changes some signs.  A cap on tuple probabilities therefore
becomes a cap on matrix entries.  We construct such a space containing
each irreducible, with dimension bounded by a fixed power of the
irreducible dimension.  This gives a normalized Hilbert--Schmidt bound.
For two half-permutations which are independent after conditioning on an
environment, a second averaging step converts that bound into an
operator-norm estimate on the full symmetric group.

For a finite set $E$, write $S_E$ for its permutation group and $U_E$
for uniform measure on that group.  An ordering of any subset $I\subseteq E$
is fixed when we write $g|_I$ or $g^{-1}|_I$ as a tuple.  The uniform
distribution on the possible tuples is independent of this choice of
ordering.  A subprobability is a nonnegative measure of total mass at
most one.  For a unitary representation $\rho_\lambda$ of dimension
$D_\lambda$, our Fourier convention is
\[
 \widehat\nu(\lambda)=\sum_g\nu(g)\rho_\lambda(g),
 \qquad \|B\|_{\mathrm{HS}}^2=\operatorname{Tr}(B^*B).
\]
Thus the Hilbert--Schmidt norm uses the ordinary, unnormalized trace.
Composition $yx$ applies $x$ first, and convolution has this same order.

\begin{theorem}[Conditional-product transfer]
\label{l:tabloid-product}
Put $C_0=2{,}000{,}000$, and fix an integer $L\ge 2C_0$.
Let $V=A\sqcup C$, where $|A|=|C|=h$ and $L$ divides $h$.
Partition each half into $L$ buffers of size $h/L$, and put
$H=S_A\times S_C\subseteq S_V$.
Let $\xi$ be a subprobability on $S_V$ such that, for $J=A,C$ and
every ordered injection $\mathbf v:J\longrightarrow V$,
\begin{equation}
 \xi\{x:x^{-1}|_J=\mathbf v\}
 \le 2U_V\{x:x^{-1}|_J=\mathbf v\}.
 \label{l:tabloid-inverse-caps}
\end{equation}
Let $Z$ have any probability distribution.  Measurably in $z$, let
$t_z\in S_V$ and let $\nu_A^z,\nu_C^z$ be subprobabilities on
$S_A,S_C$, respectively.  Suppose that for almost every $z$, each
$J=A,C$, each of its
buffers $B$, and every ordered injection
$\mathbf w:J\setminus B\longrightarrow J$,
\begin{equation}
 \nu_J^z\{y:y|_{J\setminus B}=\mathbf w\}
 \le 2U_J\{y:y|_{J\setminus B}=\mathbf w\}.
 \label{l:tabloid-factor-caps}
\end{equation}
Regard the product $\nu_A^z\otimes\nu_C^z$ as a measure on $H$ and set
\[
 \omega=\mathbb E_Z\bigl[\delta_{t_Z}*
                (\nu_A^Z\otimes\nu_C^Z)\bigr],
 \qquad \theta=\omega*\xi.
\]
Then, for all sufficiently large $h$, every irreducible representation
of $S_V$ with $D_\lambda>1$ satisfies
\begin{equation}
 \|\widehat\theta(\lambda)\|_{\mathrm{op}}
 \le K_0D_\lambda^{-1/40},\qquad K_0=\sqrt{22}.
 \label{l:tabloid-product-bound}
\end{equation}
The threshold on $h$ and the constant $K_0$ are independent of the
environment, the transporters and all the measures.
\end{theorem}

The product in the definition of $\omega$ is part of the hypothesis:
after fixing $z$, the two factors are separate subprobabilities.
The theorem allows their laws to depend arbitrarily on $z$, while the
convolution with $\xi$ uses the same measure for every environment.
The first
measure controls inverse images of whole halves; the conditional
factors control images of buffer complements.  The proof below uses
these two orientations at different steps.

\subsection{A signed tabloid containing each irreducible}

For positive part sizes $s_1,\ldots,s_r$ summing to $m$, choose on
each factor of $S_{s_1}\times\cdots\times S_{s_r}$ either its trivial
or its sign representation.  Inducing their tensor product to $S_m$
gives a signed tabloid representation.  It has an orthonormal basis
indexed by ordered set partitions $(T_1,\ldots,T_r)$ of $\{1,\ldots,m\}$
with $|T_i|=s_i$, and each group element acts by a signed permutation
matrix.  Its dimension is
\[
 M=\frac{m!}{\prod_i s_i!}.
\]
Tensoring the whole representation with sign preserves this basis
property and interchanges the trivial and sign choices on all factors.
We use the usual tableau description of irreducibles, the hook-length
formula and the Pieri rule; see \cite{sagan,etingof}.

\begin{lemma}[Polynomial-size signed tabloid embedding]
\label{l:tabloid-embedding}
For every $m\ge1$ and every partition $\lambda\vdash m$, the
irreducible representation $V_\lambda$ occurs in a signed tabloid
representation of dimension $M$ with
\[
 M\le D_\lambda^{C_0},\qquad C_0=2{,}000{,}000.
\]
\end{lemma}

\begin{proof}
The subgroup constructed in Lemma~\ref{l:small-index-character}
is a product of symmetric groups on disjoint assigned row or column
sets. Those sets partition the $m$ labels, so after omitting empty sets
it is conjugate to $J=S_{s_1}\times\cdots\times S_{s_r}$.
The occurring character $\chi$ is trivial or sign on each factor.
The full-group and trivial-subgroup choices in that lemma have the
same form.

Frobenius reciprocity now places $V_\lambda$ in
$\operatorname{Ind}_J^{S_m}\chi$. Its coset basis is indexed by the
ordered set partitions of sizes $s_1,\ldots,s_r$, and each group
element permutes this basis with signs because $\chi$ takes values
in $\{1,-1\}$. It is therefore a signed tabloid representation, with
\[
 M=[S_m:J]\le D_\lambda^{1024}\le D_\lambda^{C_0}.
\]
\end{proof}

Appendix~\ref{l:tabloid-supplement} gives an independent construction
from diagonal arms and legs, together with a degree-sum proof using
those same part sizes. The tuple estimate below needs only the signed
basis and the polynomial dimension bound just proved.

\subsection{Tuple caps give a normalized Hilbert--Schmidt bound}

\begin{lemma}[Buffer-complement estimate]
\label{l:tabloid-hs}
Let $E$ have $m$ elements, partitioned into $L$ buffers, and let
$\nu$ be a subprobability on $S_E$.  Suppose that for every buffer
$B$, the marginal of $\nu$ on the ordered image tuple of
$E\setminus B$ is pointwise at most twice its uniform marginal.
If $L\ge2C_0$, then for every irreducible $\gamma$ of $S_E$,
\begin{equation}
 \frac{\|\widehat\nu(\gamma)\|_{\mathrm{HS}}^2}{D_\gamma}
 \le \min\{1,2D_\gamma^{-1/2}\}.
 \label{l:tabloid-normalized-hs}
\end{equation}
\end{lemma}

\begin{proof}
Fix two tabloids of part sizes $s_1,\ldots,s_r$, and let $E_0$ be
the event that a permutation takes the first to the second.
For buffer $B_z$, let $s_i^z$ count the elements of the first
tabloid's part $i$ lying in $B_z$, and set
\[
 M=\frac{m!}{\prod_i s_i!},\qquad
 M_z=\frac{|B_z|!}{\prod_i s_i^z!}.
\]
Relax $E_0$ by requiring the correct target part only for positions
outside $B_z$.  This relaxed event is determined by the corresponding
image tuple.  Conditional on any successful outside assignment,
the unassigned target parts have sizes $s_i^z$, so a uniform
completion succeeds with probability $1/M_z$.
Since $U_E(E_0)=1/M$, the relaxed event has uniform probability
$M_z/M$.  Its probability under $\nu$, and hence $\nu(E_0)$,
is at most $2M_z/M$.

The product $\prod_zM_z$ counts the ways to assign the part names
inside the buffers with these fixed counts; all such assignments
are among the $M$ unrestricted assignments with totals $s_i$.
Thus $\prod_zM_z\le M$, and some buffer has $M_z\le M^{1/L}$.
We have proved the uniform entry bound
\[
 \nu(E_0)\le2M^{-1+1/L}.
\]
In the signed tabloid representation, each absolute entry of the
averaged matrix is bounded by this probability.  Each column has
absolute sum at most $\nu(S_E)\le1$, since every group element
acts by a signed permutation matrix.  Its squared Hilbert--Schmidt
norm is consequently at most $2M^{1/L}$.

Choose the signed tabloid containing $\gamma$ from
Lemma~\ref{l:tabloid-embedding}.  Orthogonal decomposition into
irreducibles shows that its squared Hilbert--Schmidt norm bounds
the contribution of any one copy of $\gamma$.  Hence
\[
 \frac{\|\widehat\nu(\gamma)\|_{\mathrm{HS}}^2}{D_\gamma}
 \le2D_\gamma^{C_0/L-1}\le2D_\gamma^{-1/2}.
\]
The bound by one follows separately because the average of
unitaries against a subprobability is a contraction.
\end{proof}

We also need to retain the multiplicity in restriction to two halves.
The following elementary consequence of the Littlewood--Richardson
rule is stronger than a count of the occurring types alone.

\begin{lemma}[Two-factor multiplicity bound]
\label{l:tabloid-lr}
For $\lambda\vdash 2h$ and $\alpha,\beta\vdash h$, let
$c_{\alpha\beta}^\lambda$ be the multiplicity of
$V_\alpha\otimes V_\beta$ in the restriction of $V_\lambda$
to $S_h\times S_h$.  Then
\[
 c_{\alpha\beta}^\lambda\le\min(D_\alpha,D_\beta).
\]
\end{lemma}

\begin{proof}
The Littlewood--Richardson rule counts certain tableaux of skew
shape $\lambda/\alpha$ and content $\beta$ whose reading words
satisfy the lattice condition.  The word determines the tableau.
Every prefix of such a word has weakly decreasing counts of the
successive letters.  All words of content $\beta$ with this prefix
condition are in bijection with standard tableaux of shape $\beta$:
put the successive integers in the next available box of the row
specified by the letter.  The prefix condition makes each partial
shape a partition.  There are $D_\beta$ such words, so the
multiplicity is at most $D_\beta$.  Interchanging the two factors
gives the bound by $D_\alpha$.
\end{proof}

\subsection{Positive conjugation and the product transfer}

We now have the two ingredients needed for
Theorem~\ref{l:tabloid-product}: a small normalized Hilbert--Schmidt
norm for each conditional factor, and a bound on the multiplicity
of each pair of factors.  The remaining step is to use the first
law's inverse-image caps to average positive operators.

\begin{proof}[Proof of Theorem~\ref{l:tabloid-product}]
Fix an irreducible $\rho=\rho_\lambda$ of $S_V$, and put
$D=D_\lambda>1$.  For almost every $z$, let
\[
 B_z=\sum_{y\in H}(\nu_A^z\otimes\nu_C^z)(y)\rho(y).
\]
Then
\[
 \widehat\theta(\lambda)
 =\mathbb E_Z\sum_x\xi(x)\rho(t_Z)B_Z\rho(x).
\]
Jensen's inequality remains valid for a subprobability by adding
the zero vector with its missing mass.  Since $\rho(t_Z)$ is
unitary, for every vector $u$ we obtain
\begin{equation}
 \|\widehat\theta(\lambda)u\|^2
 \le\mathbb E_Z\sum_x\xi(x)\|B_Z\rho(x)u\|^2.
 \label{l:tabloid-jensen}
\end{equation}
It is enough to bound the positive operator in the right side.

Fix $z$ for the moment.  The restriction to $H$ decomposes as
\[
 V_\lambda\big|_H
 =\bigoplus_{\alpha,\beta\vdash h}
      V_\alpha\otimes V_\beta\otimes
      \mathbb C^{c_{\alpha\beta}^\lambda}.
\]
On one occurring block write $a=D_\alpha$, $b=D_\beta$ and
$c=c_{\alpha\beta}^\lambda$.  Because the conditional measure is
a product, the block of $B_z$ is exactly
\[
 B_A(\alpha)\otimes B_C(\beta)\otimes I_c,
 \qquad B_J(\gamma)=\widehat{\nu_J^z}(\gamma).
\]
Both factor matrices are contractions.  The contribution of this
block to $B_z^*B_z$ is therefore bounded in positive-semidefinite
order by the operator $Q_A$ whose block is
\[
 B_A(\alpha)^*B_A(\alpha)\otimes I_b\otimes I_c
\]
and which is zero on the other blocks.  In particular $Q_A$
commutes with $\rho(S_C)$.

If $x$ and $x'$ have the same ordered inverse images on $A$, then
$x'=kx$ for some $k\in S_C$.  Thus
$\rho(x)^*Q_A\rho(x)$ depends only on $x^{-1}|_A$.
It is positive, so the pointwise cap
\eqref{l:tabloid-inverse-caps} gives an operator inequality:
\begin{align*}
 \sum_x\xi(x)\rho(x)^*Q_A\rho(x)
 &\preceq 2\mathbb E_{x\sim U_V}\rho(x)^*Q_A\rho(x)\\
 &=2\frac{\operatorname{Tr}Q_A}{D}I
 =2\frac{cab}{D}
       \frac{\|B_A(\alpha)\|_{\mathrm{HS}}^2}{a}I.
\end{align*}
Here $P\preceq Q$ means that $Q-P$ is positive semidefinite;
the uniform average is scalar by irreducibility, and its trace
determines the scalar.  Dropping the $A$ factor instead gives the
same estimate with
$\|B_C(\beta)\|_{\mathrm{HS}}^2/b$, using the cap on
$x^{-1}|_C$.  Consequently the averaged contribution of this
block is bounded by
\begin{equation}
 2\frac{cab}{D}
 \min\left\{1,2a^{-1/2},2b^{-1/2}\right\}I,
 \label{l:tabloid-one-block}
\end{equation}
where Lemma~\ref{l:tabloid-hs} supplies the two normalized norms.

The weights $cab/D$ sum to one.  For blocks with
$\max(a,b)\ge D^{1/10}$, summing
\eqref{l:tabloid-one-block} therefore costs at most
$4D^{-1/20}I$.  On each remaining block,
Lemma~\ref{l:tabloid-lr} gives $c\le D^{1/10}$, so
$cab/D\le D^{-7/10}$.  By
Theorem~\ref{l:degree-all-powers}, for all sufficiently large $h$
the number of partitions of $h$ with degree less than $D^{1/10}$
is at most
\[
 D^{1/10}\sum_{\alpha\vdash h}D_\alpha^{-1}
 \le3D^{1/10}.
\]
There are at most $9D^{1/5}$ such pairs.  Their total
contribution to \eqref{l:tabloid-one-block} is at most
$18D^{-1/2}I$.  Thus, uniformly in $z$,
\[
 \sum_x\xi(x)\rho(x)^*B_z^*B_z\rho(x)
 \preceq(4D^{-1/20}+18D^{-1/2})I
 \preceq22D^{-1/20}I.
\]
Insert this into \eqref{l:tabloid-jensen} and take square roots.
This proves \eqref{l:tabloid-product-bound}.
\end{proof}

\subsection{Mass, sign and a fixed convolution power}

\begin{corollary}[Completion of the conditional-product bound]
\label{l:tabloid-completion}
Identify $V$ with $\{1,\ldots,2h\}$, and consider a sequence of
the data of Theorem~\ref{l:tabloid-product} with $h\to\infty$
and fixed $L$.
Suppose that $\mu$ is a probability on $S_{2h}$ satisfying
$\theta\le\mu$ pointwise,
\[
 \theta(S_{2h})=1-o(1),\qquad
 \left|\widehat\mu(\mathrm{sgn})\right|=o(1).
\]
Then
\[
 \|\mu^{*100}-U_{\{1,\ldots,2h\}}\|_{\mathrm{TV}}
 \longrightarrow0.
\]
In particular the sign hypothesis holds if $\mu$ has exactly
zero expected sign.
\end{corollary}

\begin{proof}
Write $q=\theta(S_{2h})$ and $j=100$.  Convolution preserves
pointwise domination of nonnegative measures, so
$\theta^{*j}\le\mu^{*j}$, and the missing mass is $1-q^j=o(1)$.
Also
\[
 |\widehat\theta(\mathrm{sgn})|
 \le |\widehat\mu(\mathrm{sgn})|+(1-q)=o(1).
\]
The trivial coefficient of $\theta^{*j}-U$ is $q^j-1=o(1)$,
and its sign coefficient is $\widehat\theta(\mathrm{sgn})^j=o(1)$.
For all other irreducibles, operator-norm submultiplicativity
and $\|B\|_{\mathrm{HS}}\le\sqrt D\|B\|_{\mathrm{op}}$
give
\begin{align*}
 \sum_{D_\lambda>1}D_\lambda
        \|\widehat\theta(\lambda)^j\|_{\mathrm{HS}}^2
 &\le K_0^{2j}\sum_{D_\lambda>1}D_\lambda^{2-2j/40}\\
 &=K_0^{200}\sum_{D_\lambda>1}D_\lambda^{-3}=o(1),
\end{align*}
by Theorem~\ref{l:degree-all-powers}.  No normality assumption
on the Fourier matrices is used.

Finite-group Plancherel, including the trivial and sign terms,
now yields
\[
 (2h)!\sum_g\left|\theta^{*j}(g)-\frac1{(2h)!}\right|^2=o(1).
\]
Cauchy--Schwarz gives vanishing absolute sum of this difference.
Adding the nonnegative missing measure of mass $1-q^j$ proves
the asserted total-variation convergence for $\mu^{*j}$.
\end{proof}
